\documentclass{article}
\usepackage{graphicx} % Required for inserting images
\usepackage{amsmath}
\usepackage{amsfonts}
\usepackage{amssymb}
\usepackage{ytableau}
\usepackage{xcolor}
\title{Worksheet \#3}
\author{Rosa Orellana}
\date{August 5, 2026}

\begin{document}

\maketitle
\begin{enumerate}
\item Write the Schur function $s_{2,2}$ as a linear combination in the monomial, homogeneous, elementary, and power symmetric functions. 

\item Prove the following for $\lambda,\mu \vdash k$: 
\begin{enumerate}
\item If $\mu>_{\text{lex}} \lambda$, then $K_{\lambda, \mu} =0$
\item $K_{\lambda,\lambda} = 1$.
\end{enumerate}
\item Use Exercise 2 to prove that $\{s_\lambda\, : \, \lambda\vdash k\}$ is a basis for $\Lambda^k$. 

\item How many semistandard tableaux (of any shape) have content $k,k$. I.e, this means $k$ 1's and $k$ 2's.

\item Prove that $\omega$ is an involution using the identity that for $n\geq 1$, \[\sum_{i=0}^n(-1)^i e_ih_{n-i} =0.\] 

\item Use the Pieri rule to write $h_{3,2}$ in the Schur basis. 

\item Use the Pieri rule to write $e_{3,2}$ in the Schur basis. 

\item Write $p_3$ in the Schur basis. 

\item Prove or disprove: for all $\lambda, \mu\vdash k$, we have $K_{\lambda, \mu}=K_{\mu, \lambda}$.

\item For each partition of $\lambda \vdash 3$, write $h_\lambda$ as a linear combination of Schur functions. Make a conjecture about the coefficients you obtain. 

\item Use the Jacobi-Trudi identity and induction to prove that for all $n$ and $k$ with $n \geq k\geq 0$, we have 
\[ 
s_{a,1^b} = \sum_{i=0}^b (-1)^{i} h_{a+i} e_{b-i}.
\]
\end{enumerate}




\section*{Español}
\begin{enumerate}
\item Escriba la función de Schur $s_{2,2}$ como una combinación lineal de las funciones simétricas monomiales, homogéneas, elementales y de sumas de potencias.
\item Demuestre lo siguiente para $\lambda,\mu \vdash k$:
\begin{enumerate}
\item Si $\mu>_{\text{lex}} \lambda$, entonces $K_{\lambda, \mu} =0$.
\item $K_{\lambda,\lambda} = 1$.
\end{enumerate}
\item Use el Ejercicio 2 para demostrar que $\{s_\lambda\, : \, \lambda\vdash k\}$ es una base para $\Lambda^k$.
\item ¿Cuántos tableaux semiestándar (de cualquier forma) tienen contenido $k,k$? Es decir, ¿cuántos tienen $k$ entradas iguales a $1$ y $k$ entradas iguales a $2$?
\item Demuestre que $\omega$ es una involución usando la identidad \[\sum_{i=0}^n(-1)^i e_ih_{n-i} =0,\] válida para $n\geq 1$.
\item Use la regla de Pieri para escribir $h_{3,2}$ en la base de Schur.
\item Use la regla de Pieri para escribir $e_{3,2}$ en la base de Schur.
\item Escriba $p_3$ en la base de Schur.
\item Pruebe o dé un contraejemplo: para $\lambda, \mu\vdash k$, tenemos $K_{\lambda, \mu}=K_{\mu, \lambda}$.
\item Para cada partición $\lambda \vdash 3$, escriba $h_\lambda$ como una combinación lineal de funciones de Schur. Formule una conjetura sobre los coeficientes obtenidos.
\item Use la identidad de Jacobi-Trudi y la inducción para demostrar que, para todos $n$ y $k$ con $n \geq k\geq 0$, se tiene
\[
s_{a,1^b} = \sum_{i=0}^b (-1)^{i} h_{a+i} e_{b-i}.
\]
\end{enumerate}

$$
\mathbb{S}^{(2,2)}=\text{Span}\left\{ 
\begin{ytableau}
1& 2 \\
3& 4
\end{ytableau}, \quad
\begin{ytableau}
1 & 3 \\
2 & 4
\end{ytableau}
\right\}
$$
\end{document}


\begin{ytableau}
  {} & {} & {} & {} & {} & {} 
  \\  {} & {}  & {}   \\     
  {} & {}  & {}   \\
  {} & {} &*(cyan){} \\
  {}&     *(cyan){} & *(cyan){}
\end{ytableau}

\ydiagram{4,2,1}
\vskip 1in 
\begin{ytableau}
1 & 1 & 1 & 1 & ${\color{red}{2}}$ & ${\color{red}{2}}$& ${\color{blue}{2}}$ & ${\color{blue}{2}}$ & ${\color{blue}{2}}$ & ${\color{green}{3}}$ 
\\
${\color{red}{2}}$ & ${\color{blue}{2}}$  & ${\color{green}{3}}$ & ${\color{green}{3}}$ & ${\color{red}{3}}$ & ${\color{red}{3}}$ \\
${\color{red}{3}}$ 
\end{ytableau}\hskip .2in 
\vskip 1in 
\begin{ytableau}
1 & 1 & 1 & 1 & ${\color{red}{2}}$ & ${\color{red}{2}}$& ${\color{blue}{2}}$ & ${\color{green}{3}}$ & ${\color{green}{3}}$ & ${\color{green}{3}}$ 
\\
${\color{red}{2}}$ & ${\color{blue}{2}}$  & ${\color{blue}{2}}$ & ${\color{green}{3}}$ & ${\color{red}{3}}$ & ${\color{red}{3}}$ \\
${\color{red}{3}}$ 
\end{ytableau}\hskip .2in 

%\vskip 1in 

\begin{ytableau}
1 & 1  \\
2 
\end{ytableau}\hskip .2in 
\begin{ytableau}
1 & 2  \\
2 
\end{ytableau}\hskip .2in 
\begin{ytableau}
1 & 3  \\
2 
\end{ytableau}\hskip .2in 
\begin{ytableau}
1 & 1  \\
3
\end{ytableau}\hskip .2in 
\begin{ytableau}
1 & 2  \\
3 
\end{ytableau}\hskip .2in 
\begin{ytableau}
1 & 3  \\
3 
\end{ytableau}\hskip .2in 
\begin{ytableau}
2 & 2  \\
3 
\end{ytableau}\hskip .2in 
\begin{ytableau}
2 & 3  \\
3 
\end{ytableau}\hskip .2in 


\begin{ytableau}
1  \\
1  \\
1  \\
2 \\
2 \\
3
\end{ytableau}


\vskip 1in 
\begin{ytableau}
  {} & {}    \\             % First row: 3 empty boxes
  {} & {}    \\
  *(cyan){} \\
   *(cyan){} \\
    *(cyan){}
% Second row: 1 cyan box, 1 empty box       
\end{ytableau}
\vskip 1in 
\begin{ytableau}
  {} & {}  \\             % First row: 3 empty boxes
  {} & {} \\
  *(cyan){}% Second row: 1 cyan box, 1 empty box           
\end{ytableau}
